%!TEX root = ../manuscript.tex

\begin{remark}[Stacked-triple non-confluence]\label{rem:non-confluent}
The rewriting system $\mathcal{B}$ consisting of bigon R1-down and
bigon R2-down is \emph{not} confluent in general.  Consider a Gauss
word with five crossings $i,j,k,\ell,\ell'$ arranged so that
$i \succ j \succ k$ in the nesting forest, and $k$ contains two
interleaved chords $\ell,\ell'$ (the cyclic order along the boundary
is $i_{1}\,j_{1}\,k_{1}\,\ell\,\ell'\,\ell\,\ell'\,k_{2}\,j_{2}\,i_{2}$).
At this word the only available bigon moves are
$\mu_{1}=\text{R2-down}(i,j)$ and $\mu_{2}=\text{R2-down}(j,k)$.
Applying $\mu_{1}$ leaves $\{\,k,\ell,\ell'\,\}$, which is stuck
because no nested pair remains below $k$ and no adjacent pair is
present.  Applying $\mu_{2}$ leaves $\{\,i,\ell,\ell'\,\}$, which is
likewise stuck.  Thus two maximal greedy runs end in distinct
non-empty stuck states; no common descendant exists, so the system is
not confluent.

Crucially, this does \emph{not} contradict
Theorem~\ref{thm:greedy-main}, because the initial word does \emph{not}
admit any bigon reduction to $\emptyset$; both descendants are stuck.
Hence the hypothesis of the theorem is vacuously false here, and the
counterexample is perfectly consistent with
confluence-to-$\emptyset$.
\end{remark}
